% !TEX root = ../../main.tex

\section{Regularity}\label{sec:regularity}

By Proposition~\ref{thm:CLS-stationary}, the limit varifold $V$ is stationary; by Theorem~\ref{thm:integrality}, $V$ is integral in the no-mass-cancellation case, and in the mass-cancellation case the integrality of $V$ on $\partial^*\Omega(x_0) \cap \spt\|V\|$ holds unconditionally while integrality on the cancelled mass $\mu_{\mathrm{c}}$ is subject to the gap of Remark~\ref{rem:cancelled-mass-gap}. To promote $V$ to a smooth hypersurface, we apply Wickramasekera's regularity theorem (Theorem~\ref{thm:wickramasekera}), which requires two further inputs---\emph{stability} and the \emph{$\alpha$-structural hypothesis} (Definition~\ref{defn:alpha-structural}). The stability and $\alpha$-structural arguments below apply verbatim to the reduced-boundary part of $V$; under the multiplicity-one conjecture (no mass cancellation), they apply to all of $V$ and yield the full conclusion of Theorem~\ref{thm:non-excessive-main}.

%%%%%%%%%%%%%%%%%%%%%%%%%%%%%%%%%%%%%%%%%%%%%%%%%%%%%%%%%%%%%%%%%%%%%%%%%%%%%%%
\subsection{Stability}\label{subsec:stability}
%%%%%%%%%%%%%%%%%%%%%%%%%%%%%%%%%%%%%%%%%%%%%%%%%%%%%%%%%%%%%%%%%%%%%%%%%%%%%%%

The non-excessive condition provides stability through the following chain of implications: non-excessiveness ensures that the limit varifold $V$ is one-sided homotopic minimizing at all but finitely many points, and one-sided minimization implies stability.

\begin{proposition}[Stability of the limit varifold]\label{prop:stability}
Let $\Phi$ be a non-excessive ONVP sweepout, $x_i \to x_0 \in \mathfrak{m}(\Phi)$ a min-max sequence, and $V$ the varifold limit of $|\partial^*\Omega(x_i)|$. Then $V$ is stable.
\end{proposition}

\begin{proof}
By~\cite[Proposition~3.1]{CLS22}, the set $\mathfrak{h}_{\mathrm{nm}}(V)$ of non-homotopic-minimizing points (Definition~\ref{def:hnm}) is finite. Therefore $V$ is one-sided homotopic minimizing in $B_r(P)$ for every $P \in \spt\|V\| \setminus \mathfrak{h}_{\mathrm{nm}}(V)$ and some $r = r(P) > 0$. One-sided homotopic minimization implies that the second variation $\delta^2 V(\varphi, \varphi) \geq 0$ for all normal deformations $\varphi$ supported in $B_r(P)$. Since the finitely many points of $\mathfrak{h}_{\mathrm{nm}}(V)$ form a set of $\mathcal{H}^n$-measure zero, a partition-of-unity argument extends stability to all compactly supported normal deformations, giving global stability.

We remark that the finiteness of $\mathfrak{h}_{\mathrm{nm}}(V)$ in~\cite[Proposition~3.1]{CLS22} relies on both the non-excessive property and the nestedness of the sweepout: the proof constructs replacement families by ``gluing in'' one-sided homotopies, and nestedness ($\Omega(x_1) \subset \Omega(x_2)$) ensures that the parameter sets on which these replacements are defined are intervals. Nestedness is also used in the integrality argument (see Remark~\ref{rem:integrality-independence}). We also note that~\cite[Proposition~3.1]{CLS22} does not require the identification $V = |\partial^*\Omega(x_0)|$; in particular, the argument applies equally in the mass-cancellation case.
\end{proof}

It remains to verify the $\alpha$-structural hypothesis $(\mathcal{S}3)$. Note that even in the no-mass-cancellation case, where $V = |\partial^*\Omega(x_0)|$ has multiplicity one, the support $\spt\|V\|$ can have singular points where multiple branches meet (e.g., $N = 4$ half-hyperplanes in a balanced configuration); multiplicity one alone does not preclude such junctions. The following verification applies uniformly to both the mass-cancellation and no-mass-cancellation cases.

%%%%%%%%%%%%%%%%%%%%%%%%%%%%%%%%%%%%%%%%%%%%%%%%%%%%%%%%%%%%%%%%%%%%%%%%%%%%%%%
\subsection{The $\alpha$-structural hypothesis}\label{subsec:alpha-structural}
%%%%%%%%%%%%%%%%%%%%%%%%%%%%%%%%%%%%%%%%%%%%%%%%%%%%%%%%%%%%%%%%%%%%%%%%%%%%%%%

\begin{theorem}[Verification of $\alpha$-structural hypothesis]\label{thm:alpha-hypo}
    Let $V$ be the stationary integral varifold arising from a non-excessive ONVP sweepout $\Phi$
    with width $W$, via the min-max sequence $x_i \to x_0 \in \mathfrak{m}(\Phi)$. Then $V$ satisfies the $\alpha$-structural hypothesis (Definition~\ref{defn:alpha-structural}).
\end{theorem}

\begin{proof}
Suppose the $\alpha$-structural hypothesis fails at some $Z \in \sing V$. By definition, $\spt\|V\| \cap B_\rho(Z)$ consists of finitely many $C^{1,\alpha}$ hypersurfaces-with-boundary meeting along a common $C^{1,\alpha}$ boundary $\Gamma$. Blowing up at $Z$, any tangent cone $C$ consists of $N$ half-hyperplanes $P_1^+, \ldots, P_N^+$ meeting along a common edge $L$, with the stationarity balance $\sum_{j=1}^N m_j \nu_j = 0$. By Remark~\ref{rem:S-alpha-remarks}(ii), $N \geq 3$, so the pigeonhole principle gives two adjacent half-hyperplanes $P_j^+, P_{j+1}^+$ meeting at angle $\theta_j \leq 2\pi/N < \pi$.

\medskip\noindent\textbf{Step~1: Chord-beats-arc on the tangent cone.}
In each two-dimensional cross-section normal to $L$, the two half-lines $P_j^+ \cap \partial B_r$ and $P_{j+1}^+ \cap \partial B_r$ have total length $2r$, while the chord connecting their endpoints has length $2r\sin(\theta_j/2) < 2r$ (see Figure~\ref{fig:chord-beats-arc}). Integrating along $L$, the replacement surface has strictly less area:
\[
    \|C\|(B_r) - \|\widetilde{C}\|(B_r) \geq \delta(\theta_j) \cdot r^n > 0.
\]

    \begin{figure}[ht]
        \centering
        \begin{tikzpicture}[scale=1.0]
          % Edge L: shared edge of the two parallelograms, with perspective
          \coordinate (A) at (-1.8,-0.6);   % near end of edge
          \coordinate (B) at (1.8,0.6);     % far end of edge

          % H_1 (upper half-plane): parallelogram above L
          \coordinate (A1) at ($(A)+(1.2,2.4)$);
          \coordinate (B1) at ($(B)+(1.2,2.4)$);

          % H_2 (lower half-plane): parallelogram below L
          \coordinate (A2) at ($(A)+(2.4,-1.2)$);
          \coordinate (B2) at ($(B)+(2.4,-1.2)$);

          % Fill both half-planes
          \fill[fillA, opacity=0.35] (A) -- (B) -- (B1) -- (A1) -- cycle;
          \fill[fillA, opacity=0.35] (A) -- (B) -- (B2) -- (A2) -- cycle;

          % Draw H_1 edges
          \draw[thick, bindA] (A) -- (A1);
          \draw[thick, bindA] (B) -- (B1);
          \draw[bindA] (A1) -- (B1);

          % Draw H_2 edges
          \draw[thick, bindA] (A) -- (A2);
          \draw[thick, bindA] (B) -- (B2);
          \draw[bindA] (A2) -- (B2);

          % Edge L (drawn on top)
          \draw[very thick] (A) -- (B);

          % Midpoint of L and ball B_r
          \coordinate (M) at ($(A)!0.5!(B)$);
          % Circle split at tangent points: short arc (right) solid, long arc (left) dashed
          \draw[gray, thin] ($(M)+(-18:1.2)$) arc[start angle=-18, end angle=55, radius=1.2];
          \draw[gray, thin, dashed] ($(M)+(55:1.2)$) arc[start angle=55, end angle=342, radius=1.2];
          \fill[gray] (M) circle (1.5pt);
          \node[gray, above left] at ($(M)+(135:1.2)$) {$\partial B_r$};

          % Two points on L inside the circle, on either side of M
          \coordinate (P) at ($(A)!0.35!(B)$);
          \coordinate (Q) at ($(A)!0.65!(B)$);
          \fill (P) circle (1.5pt);
          \fill (Q) circle (1.5pt);

          % Blue gradient shading for chord competitor (between semi-ellipses and chord)
          \shade[left color=blue!30, right color=blue!5, shading angle=18, opacity=0.5]
            (P) .. controls (-0.11, 0.69) and (0.50, 1.14) .. (0.69, 0.99)
            -- (1.14, -0.37)
            .. controls (1.08, -0.61) and (0.33, -0.61) .. cycle;

          % Gray shading for the wedge region (between chord and arc)
          \fill[gray, opacity=0.25]
            (0.69, 0.99) -- (1.14, -0.37)
            arc[start angle=-18, end angle=55, radius=1.2]
            -- cycle;

          % Semi-ellipse on H_1 side: solid left part, dashed right part
          \draw[thick] (P) .. controls (-0.11, 0.69) and (0.50, 1.14) .. (0.69, 0.99);
          \draw[thin, dashed] (0.69, 0.99) .. controls (0.79, 0.91) and (0.77, 0.65) .. (Q);

          % Semi-ellipse on H_2 side: solid left part, dashed right part
          \draw[thick] (P) .. controls (0.33, -0.61) and (1.08, -0.61) .. (1.14, -0.37);
          \draw[thin, dashed] (1.14, -0.37) .. controls (1.18, -0.25) and (1.01, -0.05) .. (Q);

          % Chord competitor: sparse parallel blue lines clipped to competitor region
          \begin{scope}
            \clip (P) .. controls (-0.11, 0.69) and (0.50, 1.14) .. (0.69, 0.99)
              -- (1.14, -0.37)
              .. controls (1.08, -0.61) and (0.33, -0.61) .. cycle;
            \foreach \s in {-0.15, 0.1, 0.35, 0.6} {
              \draw[blue, thin] ({\s-0.9}, {\s/3+2.72}) -- ({\s+0.9}, {\s/3-2.72});
            }
          \end{scope}

          % Rightmost solid blue line: shifted slightly right of T1-T2 to cover dashed ellipses
          \draw[blue, thin] (0.72, 0.97) -- (1.16, -0.37);

          % Past tangent points: dense dashed blue lines in the wedge region
          \begin{scope}
            \clip ($(M)+(-18:1.2)$)
              arc[start angle=-18, end angle=55, radius=1.2]
              .. controls (0.79, 0.91) and (0.77, 0.65) .. (Q)
              .. controls (1.01, -0.05) and (1.18, -0.25) .. cycle;
            \foreach \s in {-0.45,-0.35,...,0.85} {
              \draw[blue, thin, dashed] ({\s-0.9}, {\s/3+2.72}) -- ({\s+0.9}, {\s/3-2.72});
            }
          \end{scope}

          % Angle arc at far end B (interior angle, the Omega_C side)
          \draw[thin] (B) ++(-27:0.7) arc[start angle=-27, end angle=63, radius=0.7];
          \node at ($(B)+(18:1.0)$) {$\theta_j$};

          % Labels
          \node[bindA, above left] at ($(A1)!0.5!(B1)$) {$P_j^+$};
          \node[bindA, below right] at ($(A2)!0.5!(B2)$) {$P_{j+1}^+$};
          \node[above left] at ($(A)!0.3!(B)$) {$L$};
        \end{tikzpicture}
        \caption{Chord-beats-arc: two adjacent half-hyperplanes $P_j^+, P_{j+1}^+$ meet along the $(n-1)$-dimensional edge $L$ at angle $\theta_j < \pi$ (by the balancing condition and pigeonhole, Remark~\ref{rem:S-alpha-remarks}(ii)). In each $2$-dimensional cross-section normal to $L$, the two half-lines (solid arcs on the shaded half-planes) have total length $2r$, while the chord (blue hatched surface) has length $2r\sin(\theta_j/2) < 2r$. The grey region is the area saved by the replacement.}
        \label{fig:chord-beats-arc}
    \end{figure}

\medskip\noindent\textbf{Step~2: Transfer to $B_r(Z)$ and the approximating sequence.}
We transfer the cone-level area deficit to the approximating Caccioppoli sets. Choose normal coordinates centered at $Z$ so that $g_{ij}(Z) = \delta_{ij}$; then $g_{ij} = \delta_{ij} + O(r)$ in $B_r(Z)$. Since the $\alpha$-structural hypothesis fails at~$Z$, the $C^{1,\alpha}$ branches of $\spt\|V\|$ in $B_\rho(Z)$ converge to the tangent cone under blow-up: at scale~$r$, each branch deviates from its limiting half-hyperplane by $O(r^{1+\alpha})$, and the angle between adjacent branches is $\theta_j + O(r^\alpha)$, which remains $< \pi$ for $r$ sufficiently small.

By varifold convergence $|\partial^*\Omega(x_i)| \to V$ and the sheet decomposition (Lemma~\ref{lem:sheet-decomposition}), for each $r \in (0, \rho)$ sufficiently small and all $i$ sufficiently large, the smooth boundary $\partial^*\Omega(x_i) \cap B_r(Z)$ decomposes into smooth sheets, each a nearly flat graph approximating one of the $C^{1,\alpha}$ branches of $\spt\|V\|$. In particular, two adjacent sheets approximate $P_j^+$ and $P_{j+1}^+$, meeting at angle $\theta_j + O(r^\alpha) < \pi$ in $B_r(Z)$.
\begin{figure}[ht]
    \centering
    \begin{tikzpicture}[scale=1.0]
      % Ball B_r(Z)
      \draw[gray, thin, dashed] (0,0) circle (2.5);
      \node[gray] at (-1.4, 2.1) {\small $\partial B_r(Z)$};

      % Tangent angle at Z for each curve
      \pgfmathsetmacro{\tA}{34}   % curve 1 tangent angle
      \pgfmathsetmacro{\tB}{-34}  % curve 2 tangent angle

      % Tangent lines at Z (gray dashed) — exact directions from \tA, \tB
      \draw[gray!60, dashed, thin] ({-2.8*cos(\tA)}, {-2.8*sin(\tA)}) -- ({2.8*cos(\tA)}, {2.8*sin(\tA)});
      \draw[gray!60, dashed, thin] ({-2.8*cos(\tB)}, {-2.8*sin(\tB)}) -- ({2.8*cos(\tB)}, {2.8*sin(\tB)});

      % Labels for P_1^+ through P_4^+ (four half-lines = rays from Z)
      \node[gray!60] at ({2.95*cos(\tB)}, {2.95*sin(\tB)-0.25}) {\small $P_1^+$};
      \node[gray!60] at ({2.95*cos(\tA)}, {2.95*sin(\tA)+0.25}) {\small $P_2^+$};
      \node[gray!60] at ({-2.95*cos(\tA)}, {-2.95*sin(\tA)+0.25}) {\small $P_3^+$};
      \node[gray!60] at ({-2.95*cos(\tB)}, {-2.95*sin(\tB)-0.25}) {\small $P_4^+$};

      % θ_1: P_1 (-34°) to P_2 (34°) — narrow angle, upper wedge
      \draw[thin] ({0.6*cos(\tB)}, {0.6*sin(\tB)}) arc[start angle=\tB, end angle=\tA, radius=0.6];
      \node at ({0.9*cos(0)}, {0.9*sin(0)}) {\small $\theta_1$};

      % θ_2: P_2 (34°) to P_3 (146°) — wide angle, left
      \draw[thin] ({0.85*cos(\tA)}, {0.85*sin(\tA)}) arc[start angle=\tA, end angle={180+\tB}, radius=0.85];
      \node at ({1.15*cos(90)}, {1.15*sin(90)}) {\small $\theta_2$};

      % θ_3: P_3 (146°) to P_4 (214°) — narrow angle, lower wedge
      \draw[thin] ({0.6*cos(180+\tB)}, {0.6*sin(180+\tB)}) arc[start angle={180+\tB}, end angle={180+\tA}, radius=0.6];
      \node at ({0.9*cos(180)}, {0.9*sin(180)}) {\small $\theta_3$};

      % θ_4: P_4 (214°) to P_1 (326°) — wide angle, right
      \draw[thin] ({0.85*cos(180+\tA)}, {0.85*sin(180+\tA)}) arc[start angle={180+\tA}, end angle={360+\tB}, radius=0.85];
      \node at ({1.15*cos(270)}, {1.15*sin(270)}) {\small $\theta_4$};

      % Upper shaded region (between curve 1 right branch and curve 2 left branch)
      \begin{scope}
      \clip (-3.2,-3.2) rectangle (3.2,3.2);
      \fill[fillA, opacity=0.2]
        (0,0)
        .. controls ({0.4*cos(\tA)}, {0.4*sin(\tA)}) and ({0.9*cos(\tA)}, {0.9*sin(\tA)}) ..
        ({2.5*cos(\tA)+0.6*sin(\tA)}, {2.5*sin(\tA)-0.6*cos(\tA)})
        -- (3.2, 3.2) -- (-3.2, 3.2)
        -- ({-2.5*cos(\tB)-0.6*sin(\tB)}, {-2.5*sin(\tB)+0.6*cos(\tB)})
        .. controls ({-0.9*cos(\tB)}, {-0.9*sin(\tB)}) and ({-0.4*cos(\tB)}, {-0.4*sin(\tB)}) ..
        (0,0);
      \end{scope}

      % Lower shaded region (between curve 2 right branch and curve 1 left branch)
      \begin{scope}
      \clip (-3.2,-3.2) rectangle (3.2,3.2);
      \fill[fillA, opacity=0.2]
        (0,0)
        .. controls ({0.4*cos(\tB)}, {0.4*sin(\tB)}) and ({0.9*cos(\tB)}, {0.9*sin(\tB)}) ..
        ({2.5*cos(\tB)+0.6*sin(\tB)}, {2.5*sin(\tB)-0.6*cos(\tB)})
        -- (3.2, -3.2) -- (-3.2, -3.2)
        -- ({-2.5*cos(\tA)-0.6*sin(\tA)}, {-2.5*sin(\tA)+0.6*cos(\tA)})
        .. controls ({-0.9*cos(\tA)}, {-0.9*sin(\tA)}) and ({-0.4*cos(\tA)}, {-0.4*sin(\tA)}) ..
        (0,0);
      \end{scope}

      % Curve 1: control points near Z lie on tangent direction \tA
      \draw[bindA, thick]
        ({-2.5*cos(\tA) - 0.6*sin(\tA)}, {-2.5*sin(\tA) + 0.6*cos(\tA)})
        .. controls ({-0.9*cos(\tA)}, {-0.9*sin(\tA)}) and ({-0.4*cos(\tA)}, {-0.4*sin(\tA)}) .. (0,0)
        .. controls ({0.4*cos(\tA)}, {0.4*sin(\tA)}) and ({0.9*cos(\tA)}, {0.9*sin(\tA)}) ..
        ({2.5*cos(\tA) + 0.6*sin(\tA)}, {2.5*sin(\tA) - 0.6*cos(\tA)});

      % Curve 2: control points near Z lie on tangent direction \tB
      \draw[bindA, thick]
        ({-2.5*cos(\tB) - 0.6*sin(\tB)}, {-2.5*sin(\tB) + 0.6*cos(\tB)})
        .. controls ({-0.9*cos(\tB)}, {-0.9*sin(\tB)}) and ({-0.4*cos(\tB)}, {-0.4*sin(\tB)}) .. (0,0)
        .. controls ({0.4*cos(\tB)}, {0.4*sin(\tB)}) and ({0.9*cos(\tB)}, {0.9*sin(\tB)}) ..
        ({2.5*cos(\tB) + 0.6*sin(\tB)}, {2.5*sin(\tB) - 0.6*cos(\tB)});

      % Blue curve (∂*Ω_i sheet) in upper wedge, just above the red boundary
      \draw[blue, thick]
        (-1.65, 2.05) .. controls (-0.8, 0.7) and (-0.2, 0.35) .. (0, 0.3)
        .. controls (0.2, 0.25) and (0.8, 0.6) .. (2.3, 1.05);

      % Blue curve (∂*Ω_i sheet) in lower wedge, rotated 180° from the upper one
      \draw[blue, thick]
        (1.65, -2.05) .. controls (0.8, -0.7) and (0.2, -0.35) .. (0, -0.3)
        .. controls (-0.2, -0.25) and (-0.8, -0.6) .. (-2.3, -1.05);

      % Point Z
      \fill (0,0) circle (1.5pt);
      \node[left] at (-0.1, 0) {$Z$};
    \end{tikzpicture}
    \caption{Cross-section normal to $L$ at $Z$: the tangent cone half-lines $P_j^+, P_{j+1}^+$ (gray dashed) and the $C^{1,\alpha}$ branches of $\spt\|V\|$ (red solid). The branches deviate from the half-lines by $O(r^{1+\alpha})$; the angle between adjacent branches is $\theta_j + O(r^\alpha) < \pi$.}
    \label{fig:transfer}
\end{figure}

\medskip\noindent\textbf{Step~3: Caccioppoli surgery and contradiction with width.}
We perform chord surgery on the approximating Caccioppoli sets. In $B_r(Z)$, the two smooth sheets of $\partial^*\Omega(x_i)$ corresponding to the narrowest adjacent pair bound a thin wedge-shaped region. Replace these two sheets inside $B_{r/2}(Z)$ by a smooth interpolation surface (the ``chord surface''), which spans the same boundary on $\partial B_{r/2}(Z)$ but cuts across the wedge.
\begin{figure}[ht]
    \centering
    \begin{tikzpicture}[scale=1.0]
      % Ball B_r(Z)
      \draw[gray, thin, dashed] (0,0) circle (2.5);
      \node[gray] at (-1.4, 2.1) {\small $\partial B_r(Z)$};

      % Tangent angle at Z for each curve
      \pgfmathsetmacro{\tA}{34}   % curve 1 tangent angle
      \pgfmathsetmacro{\tB}{-34}  % curve 2 tangent angle

      % Tangent lines at Z (gray dashed) — exact directions from \tA, \tB
      \draw[gray!60, dashed, thin] ({-2.8*cos(\tA)}, {-2.8*sin(\tA)}) -- ({2.8*cos(\tA)}, {2.8*sin(\tA)});
      \draw[gray!60, dashed, thin] ({-2.8*cos(\tB)}, {-2.8*sin(\tB)}) -- ({2.8*cos(\tB)}, {2.8*sin(\tB)});

      % Labels for P_1^+ through P_4^+ (four half-lines = rays from Z)
      \node[gray!60] at ({2.95*cos(\tB)}, {2.95*sin(\tB)-0.25}) {\small $P_1^+$};
      \node[gray!60] at ({2.95*cos(\tA)}, {2.95*sin(\tA)+0.25}) {\small $P_2^+$};
      \node[gray!60] at ({-2.95*cos(\tA)}, {-2.95*sin(\tA)+0.25}) {\small $P_3^+$};
      \node[gray!60] at ({-2.95*cos(\tB)}, {-2.95*sin(\tB)-0.25}) {\small $P_4^+$};

      % Upper shaded region (between curve 1 right branch and curve 2 left branch)
      \begin{scope}
      \clip (-3.2,-3.2) rectangle (3.2,3.2);
      \fill[fillA, opacity=0.2]
        (0,0)
        .. controls ({0.4*cos(\tA)}, {0.4*sin(\tA)}) and ({0.9*cos(\tA)}, {0.9*sin(\tA)}) ..
        ({2.5*cos(\tA)+0.6*sin(\tA)}, {2.5*sin(\tA)-0.6*cos(\tA)})
        -- (3.2, 3.2) -- (-3.2, 3.2)
        -- ({-2.5*cos(\tB)-0.6*sin(\tB)}, {-2.5*sin(\tB)+0.6*cos(\tB)})
        .. controls ({-0.9*cos(\tB)}, {-0.9*sin(\tB)}) and ({-0.4*cos(\tB)}, {-0.4*sin(\tB)}) ..
        (0,0);
      \end{scope}

      % Lower shaded region (between curve 2 right branch and curve 1 left branch)
      \begin{scope}
      \clip (-3.2,-3.2) rectangle (3.2,3.2);
      \fill[fillA, opacity=0.2]
        (0,0)
        .. controls ({0.4*cos(\tB)}, {0.4*sin(\tB)}) and ({0.9*cos(\tB)}, {0.9*sin(\tB)}) ..
        ({2.5*cos(\tB)+0.6*sin(\tB)}, {2.5*sin(\tB)-0.6*cos(\tB)})
        -- (3.2, -3.2) -- (-3.2, -3.2)
        -- ({-2.5*cos(\tA)-0.6*sin(\tA)}, {-2.5*sin(\tA)+0.6*cos(\tA)})
        .. controls ({-0.9*cos(\tA)}, {-0.9*sin(\tA)}) and ({-0.4*cos(\tA)}, {-0.4*sin(\tA)}) ..
        (0,0);
      \end{scope}

      % Red curves OUTSIDE B_{r/2} (solid — retained)
      \begin{scope}
        \pgfseteorule\clip (-3.5,-3.5) rectangle (3.5,3.5) (0,0) circle (1.25);
        \draw[bindA, thick]
          ({-2.5*cos(\tA) - 0.6*sin(\tA)}, {-2.5*sin(\tA) + 0.6*cos(\tA)})
          .. controls ({-0.9*cos(\tA)}, {-0.9*sin(\tA)}) and ({-0.4*cos(\tA)}, {-0.4*sin(\tA)}) .. (0,0)
          .. controls ({0.4*cos(\tA)}, {0.4*sin(\tA)}) and ({0.9*cos(\tA)}, {0.9*sin(\tA)}) ..
          ({2.5*cos(\tA) + 0.6*sin(\tA)}, {2.5*sin(\tA) - 0.6*cos(\tA)});
        \draw[bindA, thick]
          ({-2.5*cos(\tB) - 0.6*sin(\tB)}, {-2.5*sin(\tB) + 0.6*cos(\tB)})
          .. controls ({-0.9*cos(\tB)}, {-0.9*sin(\tB)}) and ({-0.4*cos(\tB)}, {-0.4*sin(\tB)}) .. (0,0)
          .. controls ({0.4*cos(\tB)}, {0.4*sin(\tB)}) and ({0.9*cos(\tB)}, {0.9*sin(\tB)}) ..
          ({2.5*cos(\tB) + 0.6*sin(\tB)}, {2.5*sin(\tB) - 0.6*cos(\tB)});
      \end{scope}

      % Intersection points of RED curves with ∂B_{r/2}
      % Upper wedge: right branch of curve 1 (P_2^+) and left branch of curve 2 (toward P_3^+)
      \coordinate (UR) at (1.13, 0.54);    % curve 1 right branch ∩ ∂B_{r/2}
      \coordinate (UL) at (-0.93, 0.84);   % curve 2 left branch ∩ ∂B_{r/2}
      % Lower wedge: left branch of curve 1 (P_4^+) and right branch of curve 2 (toward P_1^+)
      \coordinate (LL) at (-1.13, -0.54);  % curve 1 left branch ∩ ∂B_{r/2}
      \coordinate (LR) at (0.93, -0.84);   % curve 2 right branch ∩ ∂B_{r/2}

      % Saved area in upper wedge (between red curves and chord, inside B_{r/2})
      \begin{scope}
        \clip (0,0) circle (1.25);
        \fill[blue!12, opacity=0.6]
          (0,0)
          .. controls (0.22, 0.15) and (0.48, 0.33) .. (UR)
          -- (UL)
          .. controls (-0.48, 0.33) and (-0.22, 0.15) .. (0,0);
      \end{scope}

      % Saved area in lower wedge
      \begin{scope}
        \clip (0,0) circle (1.25);
        \fill[blue!12, opacity=0.6]
          (0,0)
          .. controls (-0.22, -0.15) and (-0.48, -0.33) .. (LL)
          -- (LR)
          .. controls (0.48, -0.33) and (0.22, -0.15) .. (0,0);
      \end{scope}

      % Blue curves (∂*Ω(x) sheets) — OUTSIDE B_{r/2} only (solid)
      \begin{scope}
        \pgfseteorule\clip (-3.5,-3.5) rectangle (3.5,3.5) (0,0) circle (1.25);
        \draw[blue, thick]
          (-1.65, 2.05) .. controls (-0.8, 0.7) and (-0.2, 0.35) .. (0, 0.3)
          .. controls (0.2, 0.25) and (0.8, 0.6) .. (2.3, 1.05);
        \draw[blue, thick]
          (1.65, -2.05) .. controls (0.8, -0.7) and (0.2, -0.35) .. (0, -0.3)
          .. controls (-0.2, -0.25) and (-0.8, -0.6) .. (-2.3, -1.05);
      \end{scope}

      % Red curves INSIDE B_{r/2} — dashed (replaced by chord surgery)
      \begin{scope}
        \clip (0,0) circle (1.25);
        \draw[bindA, thick, dashed, opacity=0.4]
          ({-2.5*cos(\tA) - 0.6*sin(\tA)}, {-2.5*sin(\tA) + 0.6*cos(\tA)})
          .. controls ({-0.9*cos(\tA)}, {-0.9*sin(\tA)}) and ({-0.4*cos(\tA)}, {-0.4*sin(\tA)}) .. (0,0)
          .. controls ({0.4*cos(\tA)}, {0.4*sin(\tA)}) and ({0.9*cos(\tA)}, {0.9*sin(\tA)}) ..
          ({2.5*cos(\tA) + 0.6*sin(\tA)}, {2.5*sin(\tA) - 0.6*cos(\tA)});
        \draw[bindA, thick, dashed, opacity=0.4]
          ({-2.5*cos(\tB) - 0.6*sin(\tB)}, {-2.5*sin(\tB) + 0.6*cos(\tB)})
          .. controls ({-0.9*cos(\tB)}, {-0.9*sin(\tB)}) and ({-0.4*cos(\tB)}, {-0.4*sin(\tB)}) .. (0,0)
          .. controls ({0.4*cos(\tB)}, {0.4*sin(\tB)}) and ({0.9*cos(\tB)}, {0.9*sin(\tB)}) ..
          ({2.5*cos(\tB) + 0.6*sin(\tB)}, {2.5*sin(\tB) - 0.6*cos(\tB)});
      \end{scope}

      % Chord surfaces S_chord (green, connecting red-curve traces on ∂B_{r/2})
      \draw[green!50!black, very thick] (UL) -- (UR);
      \draw[green!50!black, very thick] (LR) -- (LL);

      % Intersection points γ_j, γ_{j+1} on ∂B_{r/2}
      \fill (UL) circle (1.8pt);
      \fill (UR) circle (1.8pt);
      \fill (LR) circle (1.8pt);
      \fill (LL) circle (1.8pt);

      % Labels
      \node[above left, font=\small] at (UL) {$\gamma_j$};
      \node[right, font=\small] at (UR) {$\gamma_{j+1}$};
      \node[green!50!black, above, font=\small] at (0.1, 0.75) {$S_{\mathrm{chord}}$};

      % Inner ball B_{r/2}(Z)
      \draw[gray, thin, dashed] (0,0) circle (1.25);
      \node[gray] at (-0.7, -1.15) {\small $\partial B_{r/2}$};

      % Point Z
      \fill (0,0) circle (1.5pt);
      \node[left] at (-0.1, 0) {$Z$};
    \end{tikzpicture}
    \caption{Chord surgery (cross-section normal to $L$). The sheets of $\partial^*\Omega(x)$ (blue solid) approximate the $C^{1,\alpha}$ branches of $\spt\|V\|$ (red) in the narrow wedge regions. Inside $B_{r/2}(Z)$ (inner dashed circle), the original sheets (blue dashed) are replaced by chord surfaces $S_{\mathrm{chord}}$ (green), connecting the traces $\gamma_j, \gamma_{j+1}$ on $\partial B_{r/2}(Z)$. The shaded region is the area saved by the surgery: $\geq \delta(\theta_j) r^n / 2 > 0$.}
    \label{fig:surgery}
\end{figure}

\noindent\textbf{Step~3a: Chord surgery on a single slice.}
Let $\Sigma_j(x)$ and $\Sigma_{j+1}(x)$ denote the two adjacent sheets of $\partial^*\Omega(x)$ in $B_r(Z)$ approximating $P_j^+$ and $P_{j+1}^+$. Work in normal coordinates centered at $Z$ and decompose $\mathbb{R}^{n+1} = L \oplus L^\perp$, where $L$ is the $(n-1)$-dimensional edge of the tangent cone. For $\rho \in (0, r/2]$ and $y \in L$ with $|y| < \rho$, let $\Pi_y := \{y\} + L^\perp$ be the $2$-dimensional affine plane through $y$ perpendicular to $L$. Each sheet $\Sigma_k(x)$ ($k = j, j+1$) intersects the circle $\partial B_\rho(Z) \cap \Pi_y$ transversally at a point $p_k(y,x)$ on the narrow-wedge side. Define the \emph{chord surface at scale~$\rho$} as the ruled surface
\[
    S_{\mathrm{chord}}^{\rho}(x) \;:=\; \big\{(1-t)\, p_j(y,x) + t\, p_{j+1}(y,x) \,:\, y \in L \cap B_{\rho}(Z),\; t \in [0,1]\big\},
\]
i.e., in each cross-section $\Pi_y$ the two sheet-arcs inside $B_\rho(Z)$ are replaced by the straight segment connecting their trace points on $\partial B_\rho(Z)$. Let $T_{\rho}(x) \subset B_{\rho}(Z)$ denote the open region bounded by $\Sigma_j(x) \cap B_{\rho}(Z)$, $\Sigma_{j+1}(x) \cap B_{\rho}(Z)$, and $S_{\mathrm{chord}}^{\rho}(x)$---the \emph{wedge tip} cut off by the chord (see Figure~\ref{fig:surgery}). The surgered set is
\[
    \Omega(x)_\rho \;:=\; \Omega(x) \,\triangle\, T_\rho(x).
\]
Since $T_\rho(x) \subset B_\rho(Z)$, the set $\Omega(x)_\rho$ agrees with $\Omega(x)$ outside $B_\rho(Z)$, and inside $B_\rho(Z)$ the two sheets $\Sigma_j(x), \Sigma_{j+1}(x)$ are replaced by the single chord surface $S_{\mathrm{chord}}^{\rho}(x)$. The symmetric difference formulation handles both possible orientations of the wedge (whether $T_\rho(x) \subset \Omega(x)$ or $T_\rho(x) \subset M \setminus \Omega(x)$) uniformly.

In each cross-section $\Pi_y$, the chord segment $\overline{p_j\, p_{j+1}}$ has length $2\rho_y\sin(\theta_j/2)$, strictly less than the total length $2\rho_y$ of the two sheet-arcs it replaces (where $\rho_y = \sqrt{\rho^2 - |y|^2}$). Integrating along $L$ and accounting for the $C^{1,\alpha}$ deviation of the sheets and the $O(\rho)$ metric distortion:
\[
    \mathrm{Per}(\Omega(x)_\rho) \leq \mathrm{Per}(\Omega(x)) - \frac{\delta(\theta_j)}{2}\, \rho^n,
\]
where $\delta(\theta_j) > 0$ depends only on $\theta_j$ and $n$. The constant $\delta(\theta_j)$ and the radius $r$ depend on the tangent cone geometry at $Z$---the $C^{1,\alpha}$ branches of $\spt\|V\|$---but not on the individual slice; the saving is therefore uniform over all slices whose boundary decomposes into the same sheet structure.

\medskip\noindent\textbf{Step~3b: Uniformity over nearby slices.}
By continuity of $\Phi$ in the flat topology, there exists $\varepsilon > 0$ such that for every $x$ with $|x - x_0| < \varepsilon$, the boundary $\partial^*\Omega(x) \cap B_r(Z)$ admits the same sheet decomposition as described in Step~3a. The chord surgery at any scale $\rho \leq r/2$ then saves at least $\delta(\theta_j)\rho^n/2 > 0$ uniformly for all such $x$.

\medskip\noindent\textbf{Step~3c: Modified sweepout with $\mathcal{F}$-continuity.}
A na\"ive piecewise definition---performing the full surgery for $|x - x_0| < \varepsilon$ and no surgery otherwise---fails to be $\mathcal{F}$-continuous at $|x - x_0| = \varepsilon$, since the surgery modifies $\Omega(x)$ by a fixed positive volume inside $B_{r/2}(Z)$.

To ensure continuity, we use a \emph{shrinking surgery ball}. Let $\eta \colon [0,\infty) \to [0,1]$ be a smooth non-increasing function with $\eta \equiv 1$ on $[0, \tfrac{1}{2}]$ and $\eta \equiv 0$ on $[1, \infty)$. Define the \emph{shrinking surgery radius} $\rho(x) := (r/2)\cdot\eta(|x - x_0|/\varepsilon)$, so that $\rho(x) = r/2$ for $|x - x_0| \leq \varepsilon/2$ and $\rho(x) = 0$ for $|x - x_0| \geq \varepsilon$. The chord surgery at scale $\rho(x)$ is well defined whenever $\rho(x) > 0$ (the sheet decomposition holds at every scale $\leq r/2$). Set
\[
    \tilde{\Omega}(x) \;:=\; \Omega(x) \,\triangle\, T_{\rho(x)}(x),
\]
with the convention $T_0(x) := \varnothing$. Thus $\tilde{\Omega}(x) = \Omega(x)$ for $|x - x_0| \geq \varepsilon$, and $\tilde{\Omega}(x) = \Omega(x)_{r/2}$ for $|x - x_0| \leq \varepsilon/2$.

We verify that $\tilde{\Phi}(x) := \partial\tilde{\Omega}(x)$ is $\mathcal{F}$-continuous. By the triangle inequality for the flat metric,
\[
    \mathrm{Vol}\big(\tilde{\Omega}(x) \,\triangle\, \tilde{\Omega}(y)\big)
    \;\leq\; \mathrm{Vol}\big(\Omega(x) \,\triangle\, \Omega(y)\big)
    \;+\; \mathrm{Vol}\big(T_{\rho(x)}(x) \,\triangle\, T_{\rho(y)}(y)\big).
\]
The first term tends to zero as $y \to x$ by $\mathcal{F}$-continuity of $\Phi$. For the second term, note that $T_\rho(x)$ depends on $x$ through the sheets $\Sigma_j(x)$, $\Sigma_{j+1}(x)$ and the scale $\rho(x)$. As $y \to x$: the scale $\rho(y) \to \rho(x)$ by smoothness of $\eta$; and the sheets converge in $C^1(B_r(Z))$, since flat convergence $\Omega(y) \to \Omega(x)$ combined with the uniform perimeter bound (all slices have $\mathrm{Per} \leq W + 1$) implies varifold convergence of boundaries, which for graphical sheets with small $C^{1,\alpha}$ norm over the tangent half-hyperplanes upgrades to smooth convergence by standard graphical estimates~\cite[Theorem~24.2]{Sim84}. Therefore $\mathrm{Vol}(T_{\rho(x)}(x) \,\triangle\, T_{\rho(y)}(y)) \to 0$, and $\tilde{\Phi}$ is $\mathcal{F}$-continuous on $[0,1]$.

\medskip\noindent\textbf{Step~3d: Replacement family and contradiction.}
Since $\theta_j < \pi$, the chord surface has strictly less area than the two sheets it replaces at every scale: for any $x$ with $|x - x_0| < \varepsilon$,
\[
    \mathrm{Per}\big(\tilde{\Omega}(x)\big) \leq \mathrm{Per}\big(\Omega(x)\big) - \frac{\delta(\theta_j)}{2}\,\rho(x)^n.
\]
Set $I := (x_0 - \varepsilon,\, x_0 + \varepsilon)$. Since $\tilde{\Omega}(x) = \Omega(x)$ at the endpoints $x = x_0 \pm \varepsilon$, we claim that $\tilde{\Phi}|_{\bar{I}}$ is an $I$-replacement family for $\Phi$ (Definition~\ref{def:excessive-interval}). Indeed, for each $x \in I$ we have $\rho(x) > 0$ and $\rho$ is continuous, so
\begin{align*}
    \limsup_{I \ni y \to x}\, \mathrm{Per}\big(\tilde{\Omega}(y)\big)
    &\;\leq\; \limsup_{y \to x}\, \mathrm{Per}\big(\Omega(y)\big)
    - \frac{\delta(\theta_j)}{2}\,\rho(x)^n \\
    &\;\leq\; W - \frac{\delta(\theta_j)}{2}\,\rho(x)^n
    \;<\; W,
\end{align*}
where the second inequality uses $\limsup_{y \to x}\,\mathbf{M}(\Phi(y)) = M(x) \leq W$.
Therefore $I$ is an excessive interval for $\Phi$ containing $x_0$, which means $x_0$ is excessive---contradicting the non-excessive property of $\Phi$ (Definition~\ref{def:excessive-interval}).

\medskip
Therefore, the $\alpha$-structural hypothesis holds.
\end{proof}

We can now assemble the full regularity result.

\begin{proof}[Proof of Theorem~\ref{thm:non-excessive-main}]
Let $(M^{n+1}, g)$ be a closed Riemannian manifold with $n \geq 2$. By Theorem~\ref{thm:non-excessive-existence}, there exists a non-excessive ONVP sweepout $\Phi \in \Snm$. By Proposition~\ref{thm:CLS-stationary}, there exists a stationary varifold $V$ with $\mathbf{M}(V) = W$.

By Theorem~\ref{thm:integrality}(a), $V$ is an integral varifold in the no-mass-cancellation case (and, conditionally on Remark~\ref{rem:cancelled-mass-gap}, also in the mass-cancellation case). By Proposition~\ref{prop:stability}, $V$ is stable. By Theorem~\ref{thm:alpha-hypo}, $V$ satisfies the $\alpha$-structural hypothesis. Therefore $V \in \mathcal{S}_\alpha$, and Wickramasekera's theorem (Theorem~\ref{thm:wickramasekera}) gives the stated regularity: $\sing V = \varnothing$ for $2 \leq n \leq 6$, $\sing V$ discrete for $n = 7$, and $\mathcal{H}^{n-7+\gamma}(\sing V) = 0$ for each $\gamma > 0$ when $n \geq 8$. Thus Theorem~\ref{thm:non-excessive-main} holds unconditionally under no mass cancellation, and otherwise holds conditionally on the resolution of Remark~\ref{rem:cancelled-mass-gap}.
\end{proof}
